\documentclass[aspectratio=169,10pt]{beamer}
\usetheme{Madrid}
\usecolortheme{dolphin}
\usepackage[utf8]{inputenc}
\usepackage[T1]{fontenc}
\usepackage[portuguese]{babel}
\usepackage{lmodern}
\usepackage{amsmath}
\usepackage{booktabs}
\usepackage{tikz}
\usetikzlibrary{arrows.meta,positioning}
\usepackage{listings}
\definecolor{azul}{RGB}{18,78,133}
\definecolor{verde}{RGB}{26,128,91}
\setbeamercolor{structure}{fg=azul}
\setbeamercolor{frametitle}{bg=azul,fg=white}
\setbeamercolor{block title}{bg=azul,fg=white}
\setbeamertemplate{navigation symbols}{}
\setbeamertemplate{footline}[frame number]
\setbeamersize{text margin left=8mm,text margin right=8mm}
\lstset{language=C,basicstyle=\ttfamily\scriptsize,keywordstyle=\color{azul}\bfseries,showstringspaces=false,columns=fullflexible,keepspaces=true,frame=single,framerule=0.3pt,xleftmargin=3pt,xrightmargin=3pt}
\newcommand{\etiqueta}[1]{{\color{azul}\bfseries #1}}
\begin{document}

% SLIDE 1
\begin{frame}{Sistema de aquisicao e monitorizacao de vibracoes}
\begin{columns}[T,onlytextwidth]
\column{0.40\textwidth}
{\Large\bfseries Trabalho Final --- SISEM}\\[3mm]
\etiqueta{Objetivo}\\
Medir e processar a intensidade de vibracoes nos eixos X, Y e Z.\\[4mm]
\etiqueta{Hardware}\\
STM32F103RBT6\\
Acelerometro ADXL337\\
LCD integrado\\[4mm]
\textbf{Grupo:} [Nomes dos alunos]
\column{0.60\textwidth}
\centering
\begin{tikzpicture}[>=Stealth, every node/.style={font=\scriptsize},box/.style={draw=azul,rounded corners=2pt,fill=blue!5,align=center,minimum width=21mm,minimum height=8mm}]
\node[box] (sensor) at (0,2.45) {ADXL337\\X, Y, Z};
\node[box] (adc) at (0,1.35) {ADC1 + TIM3};
\node[box] (dma) at (0,0.25) {DMA1 + RAM};
\node[box] (rms) at (0,-0.85) {RMS + Med5};
\node[box] (uart) at (-1.8,-2.1) {USART2\\Computador};
\node[box] (lcd) at (1.8,-2.1) {SPI1 + LCD\\Barra / balao};
\draw[->,thick] (sensor)--(adc);
\draw[->,thick] (adc)--(dma);
\draw[->,thick] (dma)--(rms);
\draw[->,thick] (rms)--(uart);
\draw[->,thick] (rms)--(lcd);
\end{tikzpicture}
\end{columns}
\end{frame}

% SLIDE 2
\begin{frame}{ADXL337: sensor e ligacoes ao ADC}
\begin{columns}[T,onlytextwidth]
\column{0.52\textwidth}
\begin{block}{Caracteristicas essenciais}
\begin{itemize}\setlength\itemsep{2pt}
\item Acelerometro \textbf{MEMS de 3 eixos}.
\item Saidas \textbf{analogicas} XOUT, YOUT e ZOUT.
\item Gama de medicao de pelo menos $\pm3\,g$.
\item Sensibilidade tipica: \textbf{300 mV/g a 3,0 V}.
\item Para 3,3 V: estimativa nominal de \textbf{330 mV/g}.
\end{itemize}
\end{block}
\column{0.45\textwidth}
\begin{block}{Ligacoes utilizadas}
\centering\begin{tabular}{lll}
\toprule
Eixo & Pino & Canal ADC1\\\midrule
X & PC2 & IN12\\
Y & PC1 & IN11\\
Z & PC0 & IN10\\\bottomrule
\end{tabular}
\end{block}
\vspace{2mm}
\centering
\fbox{\parbox{0.90\linewidth}{\centering \textbf{ADC com 12 bits}\\[1mm]
$0\ \ldots\ 4095$\\[1mm]
\scriptsize Digitaliza a tensao de cada saida}}
\end{columns}
\vfill
{\footnotesize\color{gray}Fonte tecnica: datasheet ADXL337 (Analog Devices). A conversao em mg e nominal, nao uma calibracao de precisao.}
\end{frame}

% SLIDE 3
\begin{frame}[fragile]{ADC1 + TIM3: aquisicao periodica}
\begin{columns}[T,onlytextwidth]
\column{0.50\textwidth}
\etiqueta{Decisoes de configuracao}
\begin{itemize}\setlength\itemsep{3pt}
\item \textbf{Scan}: 3 canais na ordem X $\to$ Y $\to$ Z.
\item \textbf{Continuous = DISABLE}: cada sequencia comeca por disparo externo.
\item \textbf{TIM3 TRGO}: frequencia regular, independente da interface grafica.
\item Clock ADC = \textbf{12 MHz}.
\end{itemize}
\vspace{2mm}
\[f_s=\frac{72\,000\,000}{(7199+1)(49+1)}=\boxed{200\ \mathrm{Hz}}\]
\centering\small Uma sequencia de 3 conversoes a cada \textbf{5 ms}.
\column{0.48\textwidth}
\etiqueta{Exemplo de configuracao SPL}
\begin{lstlisting}
ADC_InitStructure.ADC_ScanConvMode
    = ENABLE;
ADC_InitStructure.ADC_ContinuousConvMode
    = DISABLE;
ADC_InitStructure.ADC_ExternalTrigConv
    = ADC_ExternalTrigConv_T3_TRGO;
ADC_InitStructure.ADC_NbrOfChannel = 3;
ADC_Init(ADC1, &ADC_InitStructure);
\end{lstlisting}
\vspace{3mm}
\centering
\begin{tikzpicture}[>=Stealth,font=\scriptsize]
\draw[->] (0,0)--(5.0,0) node[right]{tempo};
\draw[azul,thick] (.5,-.2)--(.5,.4);\node[above] at (.5,.4) {X Y Z};
\draw[azul,thick] (2.5,-.2)--(2.5,.4);\node[above] at (2.5,.4) {X Y Z};
\draw[azul,thick] (4.5,-.2)--(4.5,.4);\node[above] at (4.5,.4) {X Y Z};
\draw[<->] (.5,-.45)--(2.5,-.45) node[midway,below]{5 ms};
\end{tikzpicture}
\end{columns}
\end{frame}

% SLIDE 4
\begin{frame}[fragile]{DMA circular: adquirir sem bloquear a CPU}
\begin{columns}[T,onlytextwidth]
\column{0.53\textwidth}
\begin{block}{Como funciona}
\begin{itemize}\setlength\itemsep{2pt}
\item DMA1, \textbf{canal 1}: $\mathrm{ADC1\!\to RAM}$.
\item Transferencias de \textbf{16 bits}; endereco RAM incrementado.
\item Buffer \textbf{circular} com duas metades.
\item Flags \textbf{HT / TC}: sinalizam uma metade concluida.
\item Copia local da metade antes do calculo RMS.
\end{itemize}
\end{block}
\begin{lstlisting}
volatile uint16_t adc_buffer[2][32][3];
DMA_InitStructure.DMA_Mode =
    DMA_Mode_Circular;
DMA_Init(DMA1_Channel1, &DMA_InitStructure);
\end{lstlisting}
\column{0.44\textwidth}
\centering\etiqueta{Organizacao do buffer}\\[3mm]
\begin{tikzpicture}[font=\small]
\draw[thick,azul,fill=blue!8] (0,0) rectangle (4,1);
\draw[thick,azul,fill=green!10] (0,-1) rectangle (4,0);
\node at (2,.5) {Metade 0: 32 $\times$ 3};
\node at (2,-.5) {Metade 1: 32 $\times$ 3};
\draw[->,thick,azul] (4.15,.5) arc[start angle=80,end angle=-80,radius=.9];
\end{tikzpicture}
\vspace{5mm}
\[T_{\text{bloco}}=\frac{32}{200}=\boxed{160\ \mathrm{ms}}\]
\small Uma metade e processada enquanto a aquisicao continua.
\end{columns}
\end{frame}

% SLIDE 5
\begin{frame}{Processamento: RMS, mg e mediana de 5 blocos}
\begin{columns}[T,onlytextwidth]
\column{0.53\textwidth}
\etiqueta{RMS da componente variavel}\\[-1mm]
\[RMS_X=\sqrt{\frac{1}{32}\sum_{i=1}^{32}(x_i-\bar x)^2}\]
\[RMS_{\mathrm{total}}=\sqrt{RMS_X^2+RMS_Y^2+RMS_Z^2}\]
\small
Subtrair $\bar x$ remove a componente constante de cada janela (incluindo parte do efeito da gravidade).\\[2mm]
\etiqueta{Conversao nominal}\\
\[1g\approx409{,}5\ \text{contagens ADC}\]
\column{0.44\textwidth}
\begin{block}{Por que usamos a Med5?}
\begin{itemize}\setlength\itemsep{3pt}
\item Guardamos os ultimos \textbf{5 RMS totais}.
\item Ordenamos e escolhemos o \textbf{valor central}.
\item Atenua picos e quedas \textbf{isolados}.
\item Controla a animacao sem alterar o RMS original.
\end{itemize}
\end{block}
\vspace{1mm}
\small\textbf{Exemplo:} $20,22,85,21,19$ mg\\
Mediana: \textbf{21 mg}; media: 33,4 mg.\\[3mm]
\textbf{USART2 (9600 bps):} X, Y, Z, Total e Med5.
\end{columns}
\vfill
{\scriptsize\color{gray}A Med5 nao elimina a causa da interferencia do LCD nem recupera informacao perdida por aliasing.}
\end{frame}

% SLIDE 6
\begin{frame}{Demonstracao no LCD e conclusoes}
\begin{columns}[T,onlytextwidth]
\column{0.52\textwidth}
\etiqueta{Aplicacao de monitorizacao}
\begin{itemize}\setlength\itemsep{3pt}
\item Barra e balao crescem com a intensidade filtrada.
\item Em repouso, zona morta de \textbf{35 mg}.
\item Niveis \textbf{BAIXA / MEDIA / ALTA} com histerese.
\item Rebentamento associado a \textbf{Med5 $\geq$ 500 mg}, com confirmacao; reinicio automatico.
\item Desenho incremental e atualizacoes espacadas para reduzir atividade SPI.
\end{itemize}
\column{0.45\textwidth}
\centering\etiqueta{Esquema da interface}\\[3mm]
\begin{tikzpicture}[scale=.82]
\draw[azul,very thick,rounded corners=4pt,fill=black!3] (0,0) rectangle (4.8,4.6);
\node[font=\small\bfseries] at (2.4,4.12) {SISEM VIBRACAO};
\node[font=\small,text=verde] at (2.4,3.55) {VIB: MEDIA};
\draw[fill=green!10] (.65,2.85) rectangle (4.15,3.14);
\fill[verde] (.65,2.85) rectangle (2.8,3.14);
\draw[fill=red!75,draw=red!75] (2.4,1.55) circle (.70);
\draw[gray,thick] (2.4,.85)--(2.4,.25);
\end{tikzpicture}
\end{columns}
\vfill
\begin{block}{Conclusao}
Implementamos a cadeia completa: \textbf{sensor $\to$ ADC + TIM3 $\to$ DMA $\to$ RMS/Med5 $\to$ USART + LCD}. A interface e demonstrativa; os limites em mg nao sao limiares de seguranca calibrados.
\end{block}
\end{frame}

\end{document}
